\section{Valence six: a ladder and a factorial}\label{sec:v6}

At valence six the sign configurations are too many for a normal form (there are
$23\,808$ of them~\cite[Table~1]{BR2008b}), and we treat only the configuration that
Brunnemann and Rideout single out as producing the smallest eigenvalues,
\[
  \sigma=(2,0,0,0,0,0,0,0,-2,0)\quad\text{in the order }123,124,125,134,135,145,234,235,245,345,
\]
so that $A=2\hat q_{123}-2\hat q_{245}=-8i(X-Y)$ with $X=J_1\cdot(J_2\times J_3)$ and
$Y=J_2\cdot(J_4\times J_5)$. It is realised by the integer tangents
$\dot e_1=(-1,2,-3)$, $\dot e_2=(-2,4,1)$, $\dot e_3=(2,-1,-4)$, $\dot e_4=(-3,1,4)$,
$\dot e_5=(-4,3,3)$, $\dot e_6=(-2,-3,2)$, all of whose twenty triple determinants are
nonzero. The spins are $2j=(1,1,n,n,n+1,n+1)$, Brunnemann and Rideout's minimising
spins at even $2j_{\max}=n+1$~\cite[Table~9]{BR2008b}; at odd $2j_{\max}$ their minimisers
are different. Their Table~9 labels the configuration both with $-2$ and with $+2$ in the
ninth place; both labels are realisable in general position, and they give the same ladder and
identical spectra.

\subsection{The picture: two rows of dimers}

The two triple products share leg $2$, so the useful labels are the spin
$c=|J_1+J_3|\in\{c_-,c_+\}=\{\frac{n-1}2,\frac{n+1}2\}$ of legs $1,3$ and the spin
$e=|J_4+J_5|$ of legs $4,5$. In the basis that also fixes $h=|J_2+J_4+J_5|$ the
product $Y$ is a set of dimers: it pairs $e=h-\frac12$ with $e=h+\frac12$ on each row
$c$, with weight $y_h=\frac14\sqrt{(n+1-h)(n+2+h)h(h+1)}$, and it has exactly two zero
rows, at the two ends of the row $c_+$. The product $X$ is a rung: in each block of fixed
$e$ it is a rank-two antisymmetric map of fixed size $x=\frac14\sqrt{n(n+2)}$ between
the two rows. So the matrix is a two-row ladder (\cref{fig:ladder}). The two zero rows
of $Y$ are end modes, and to couple them through the ladder costs a factor $x/y_h$ per
dimer, about $1/h$ for $h$ small against $n$. After $n$ dimers the cost is the product,
and it can be written down exactly: $\prod_{h=1}^n(x/y_h)^2$ is $(n(n+2))^n$ divided by
$n!\,(n+1)!\,n!\,(2n+2)!/(n+2)!$, which by Stirling is $\exp(-2n\ln n+O(n))$. So the cost
is about $1/n!$, up to a factor $e^{O(n)}$.

\begin{schemfigure}
\begin{tikzpicture}[scale=0.74]
  \foreach \i in {0,...,7}{
    \draw[esPlotAmber, line width=2.2pt] (2*\i+0.25,1.6) -- (2*\i+1.05,1.6);
    \draw[esPlotAmber, line width=2.2pt] (2*\i+0.25,0) -- (2*\i+1.05,0);
    \fill[esPlotInk] (2*\i+0.25,1.6) circle (2.6pt); \fill[esPlotInk] (2*\i+1.05,1.6) circle (2.6pt);
    \fill[esPlotInk] (2*\i+0.25,0) circle (2.6pt); \fill[esPlotInk] (2*\i+1.05,0) circle (2.6pt);
    \draw[esPlotBlue, line width=0.6pt] (2*\i+1.05,1.6) -- (2*\i+1.65,0);
    \draw[esPlotBlue, line width=0.6pt] (2*\i+1.05,0) -- (2*\i+1.65,1.6);
  }
  \draw[esPlotInk, fill=white, line width=1pt] (-0.55,1.6) circle (3pt);
  \draw[esPlotInk, fill=white, line width=1pt] (16.05,1.6) circle (3pt);
  \draw[esPlotBlue, line width=0.6pt] (-0.55,1.6) -- (0.25,0);
  \draw[esPlotBlue, line width=0.6pt] (15.05,0) -- (16.05,1.6);
  \node[anchor=east, font=\small] at (-0.9,1.6) {row $c_+$};
  \node[anchor=east, font=\small] at (-0.9,0) {row $c_-$};
  \node[font=\small, text=esPlotAmber!80!black] at (4.65,2.15) {$Y$: dimers $y_h\approx\frac14 nh$};
  \node[font=\small, text=esPlotBlue] at (10.6,2.15) {$X$: rungs of size $x\approx\frac14 n$};
  \node[font=\small, align=center] at (-0.55,2.55) {end mode};
  \node[font=\small, align=center] at (16.05,2.55) {end mode};
\end{tikzpicture}
\end{schemfigure}
\esfiglabel{fig:ladder}%
\schemcap{FIG.\ \theesfig}{The six-valent operator as a ladder. Each row is a chain of
dimers of $Y$; the rungs of $X$ couple the rows only through a fixed direction in each
block. The two sites that $Y$ leaves unpaired, at the ends of the upper row, carry the
nearly-zero mode; it reaches across through $n$ dimers, each costing a factor of
order $x/y_h\approx1/h$, so the splitting is factorial, not exponential.}

\subsection{The theorem}

\begin{theorem}{Factorial gap closing at valence six}{v6}
For the configuration and spins above:
\begin{enumerate}
\item[(i)] for odd $n$ the operator has trivial kernel, and for even $n$ its kernel is
  exactly two-dimensional;
\item[(ii)] for odd $n\ge7$,
  \[
    0<\qm(n)\ \le\ 22700\Bigl(\frac{(n+1)\sqrt{2n}}4+\frac{\sqrt n}4\Bigr)\sqrt{12}\;
    \frac{(n+2)^{(n-3)/2}}{((n-2)!)^{3/2}}\ =\ \exp\bigl(-n\ln n+O(n)\bigr).
  \]
\end{enumerate}
So the smallest nonzero volume tends to zero along a six-valent embedding in general position,
faster than any exponential.
\end{theorem}

The proof is in \cref{app:valence6}. It runs the zero-energy recursion along the ladder:
the equations of $A\psi=0$ split into two independent chains of $2\times2$ transfer
steps, one killed by the left end, the other carrying the end modes, which reaches the
right end only for odd $n$ (this is (i)). For odd $n$ a quasimode built from the
transfer matrix $\mathsf T$ gives $\smin\le\mathrm{const}\cdot|\det\mathsf T|/\|\mathsf T\|$; the determinant
telescopes to a bounded quantity, and $\|\mathsf T\|$ grows like the product $\prod(y_j/x)$ over
the dimers, whose square is $(n-1)!\,n!\,(n-2)!\,\prod_{k=n+5}^{2n+1}k/(12(n(n+2))^{n-3})$.
The proof uses the moduli of the matrix elements and the sign structure of
\cref{prop:ladder}, in which each rung block has rank two; it does not cover arbitrary
independent signs of the entries.

\paragraph{The rate, and Brunnemann and Rideout's fit.}
By Stirling, $\ln\prod_{j=3}^{n-1}(y_j/x)=n\ln n+O(n)$, so the bound is
$\exp(-n\ln n+O(n))$. Exact characteristic polynomials of the operator for odd
$n\le51$ (computed from \eqref{eq:sigma}, and identical to those of the full
twenty-triple operator for $n\le6$) give $\ln\qm=-1.0016\,n\ln n+0.316\,n+\dots$ for a four-parameter fit over
$15\le n\le51$, consistent with the bound having the right leading term. (The factor
$((n-2)!)^{3/2}$ in the bound alone would suggest a coefficient $\frac32$; the numerator
$(n+2)^{(n-3)/2}$ contributes $+\frac12n\ln n$, and the coefficient is $1$.) In Brunnemann and Rideout's variables
this is $\lambda_V=\sqrt{\qm}\approx\exp(-\frac12n\ln n)$, a factorial law; their fit
$\lambda_V\approx e^{3.72}e^{-1.215\cdot2j_{\max}}$ over $2j_{\max}\le13$ is an
exponential fit to it over a short window, where $n\ln n$ is close to linear, and it
should not be extrapolated. Our own values on this family are $\qm=0.840877$,
$0.0459024$, $6.86354\times10^{-4}$ and $4.65514\times10^{-6}$ for $n=1,3,5,7$.

\begin{remark}[Every valence, degenerately]
If several edges share one tangent direction, every triple containing two of them has
$\epsilon=0$, and the operator of an $N$-valent vertex is that of a five-valent vertex
with one leg replaced by the total spin of the bundle, because the triple products of
the bundle's edges with any two other legs carry the same sign and add up to the triple
product of the bundle's total spin, while those with two edges inside the bundle have
$\epsilon=0$. So bundling edges onto the fifth edge of the vertex of
\cref{thm:closing} closes the gap at every valence $N\ge5$. These vertices are five-valent
in disguise, and we do not count them as evidence about generic higher valence.
\end{remark}
